
% Marco teórico adaptado desde el borrador proporcionado por el estudiante.
% Las marcas \pendiente deben resolverse antes de la versión final.
\newcommand{\pendiente}[1]{\textcolor{red!70!black}{\textsf{[pendiente: #1]}}}

% Notación empleada en este capítulo.
\newcommand{\dom}{\Omega}
\newcommand{\pde}{\mathcal{N}}
\newcommand{\res}{r}
\newcommand{\sig}{s}
\newcommand{\pesof}{w}
\newcommand{\despl}{\bm{d}}
\newcommand{\potel}{\varphi}
\newcommand{\univ}{\psi}

\label{sec:marco-teorico}


\section{El objeto y su dificultad}

La aproximación numérica de ecuaciones diferenciales parciales sostiene buena parte de la ingeniería y las ciencias aplicadas. El problema puede enunciarse en forma general,  dado un operador diferencial $\pde$ sobre un dominio $\dom$ y un conjunto de condiciones sobre $\partial\dom$, se busca $u$ tal que
\begin{equation}
  \pde[u(x)] = 0, \quad x \in \dom, \qquad
  B[u(x)] = 0, \quad x \in \partial\dom .
\label{eq:pde-general}
\end{equation}


La dificultad de resolver \eqref{eq:pde-general} no se distribuye de manera uniforme sobre $\dom$. En la mayoría de los problemas de interés , la solución exige proporcionalmente más capacidad de aproximación en regiones acotadas del dominio que en el resto, capas límite, frentes de onda, interfaces entre materiales, discontinuidades geométricas y zonas de concentración de esfuerzo.

\pendiente{Aquí entra la evidencia concreta de Daniel. \cite{gonzalez2026pinnpiezo}: los mapas de error puntual del efecto piezoeléctrico indirecto localizan la discrepancia en la interfaz entre capas de polarización opuesta. Comprimir y remitir a la sección del caso piezoeléctrico para el desarrollo.}


\section{Precedente: refinamiento donde la solución lo exige}

La concentración local de dificultad no es un descubrimiento reciente ni una particularidad de los métodos basados en aprendizaje. Los métodos numéricos clásicos la reconocieron y construyeron mecanismos explícitos para responder a ella. El refinamiento adaptativo de malla asigna más elementos donde un estimador local de error lo indica, mientras que el refinamiento $hp$ combina la reducción del tamaño de los elementos con el aumento del orden polinomial en las regiones críticas.

\pendiente{Agregar dos o tres oraciones sobre estimadores de error a posteriori, con cita. No desarrollar la teoría: la sección existe para establecer que redistribuir grados de libertad es una necesidad reconocida, no una ocurrencia de este trabajo.}

El método de elementos finitos conserva ventajas de precisión y de garantía teórica que ningún método basado en redes ha igualado \cite{lu2021deepxde}. Su costo, sin embargo, crece con la complejidad de la malla, y las formulaciones inversas o los escenarios que requieren inferencia rápida lo vuelven poco práctico \cite{karniadakis2021}. Esa tensión, y no una supuesta obsolescencia del método clásico, es la que abre espacio a las arquitecturas informadas por física.

\section{Redes informadas por física y la capacidad que se pierde}

Las redes neuronales informadas por física aproximan la solución de \eqref{eq:pde-general} con una red $u_\theta$ e incorporan la información física en la función de pérdida \cite{raissi2019}. El residual se define como la aplicación del operador diferencial a la aproximación,
\begin{equation}
  \res(x;\theta) := \pde[u_\theta(x)],
\label{eq:residual}
\end{equation}

y la pérdida combina el residual evaluado en un conjunto de puntos de colocación del interior con un término que penaliza el incumplimiento de las condiciones de frontera:
\begin{equation}
  \mathcal{L}(\theta)
  = \frac{1}{N}\sum_{i=1}^{N} \left\| \res(x_i;\theta) \right\|^2
  + \frac{1}{M}\sum_{j=1}^{M} \left\| B[u_\theta(x_j)] \right\|^2 .
\label{eq:perdida}
\end{equation}


La formulación prescinde de la malla y admite con naturalidad la estimación de parámetros desconocidos a partir de observaciones parciales, que es donde el enfoque ofrece ventaja frente a los métodos clásicos \cite{raissi2019}.

Un perceptrón multicapa parametriza la solución de forma global,  cada peso influye sobre el valor de $u_\theta$ en todo el dominio, y no existe un subconjunto de parámetros asociado a una región. Los puntos de colocación en \eqref{eq:perdida} se distribuyen además de manera uniforme o cuasiuniforme, sin relación con dónde la solución exige más capacidad. El resultado es que el mecanismo de refinamiento local que la numérica clásica había construido desaparece del método.

El sesgo espectral es la manifestación medible de esa pérdida, las redes ajustan primero las componentes de baja frecuencia de la solución y convergen con lentitud sobre las de alta frecuencia
\pendiente{Agregar y verificar una cita: Rahaman et al. (2019) o Wang, Yu y Perdikaris (2022) para la versión específica de PINNs vía núcleo tangente neuronal.}
En problemas con soluciones oscilatorias, comportamiento multiescala o acoplamiento físico fuerte, la degradación se concentra  en las regiones que determinan la validez física de la solución.

\section{Estrategias de reparación existentes}

La literatura respondió a esta pérdida con estrategias que redistribuyen el esfuerzo del entrenamiento. El refinamiento adaptativo basado en residuales añade puntos de colocación donde el residual es mayor \cite{lu2021deepxde}. Las variantes de muestreo por importancia extienden esa idea con distribuciones de probabilidad sobre el dominio.
\pendiente{Agregar una cita concreta de muestreo adaptativo; candidatos en el mapa de referencias del proyecto.}
Las PINNs variacionales con refinamiento $hp$ adaptan este precedente al paradigma de aprendizaje informado por física mediante la partición del dominio en subdominios y el control del orden de las funciones de prueba \cite{kharazmi2021hpvpinns}. El entrenamiento causal introduce pesos que penalizan los residuales de tiempos posteriores antes de resolver los anteriores \cite{wang2024causal}.

\begin{table}[htbp]
  \centering
  \small
  \caption{Estrategias existentes sobre los cuatro ejes de evaluación del capítulo.}
  \label{tab:rejilla}
  \resizebox{\textwidth}{!}{%
    \begin{tabular}{lcccc}
      \toprule
      Estrategia & Error local & Residual & Estabilidad & Costo \\
      \midrule
      RAR \cite{lu2021deepxde} & \textit{Pendiente} & \textit{Pendiente} & \textit{Pendiente} & \textit{Pendiente} \\
      Muestreo por importancia & \textit{Pendiente} & \textit{Pendiente} & \textit{Pendiente} & \textit{Pendiente} \\
      $hp$-VPINN \cite{kharazmi2021hpvpinns} & \textit{Pendiente} & \textit{Pendiente} & \textit{Pendiente} & \textit{Pendiente} \\
      Entrenamiento causal \cite{wang2024causal} & \textit{Pendiente} & \textit{Pendiente} & \textit{Pendiente} & \textit{Pendiente} \\
      Adaptación de grilla KAN (fuente nueva por verificar) & \textit{Pendiente} & \textit{Pendiente} & \textit{Pendiente} & \textit{Pendiente} \\
      \bottomrule
    \end{tabular}%
  }
\end{table}


Las estrategias comparten una limitación de estructura al reasignar dónde se evalúa el residual o con qué peso se pondera cada término de la pérdida. Ninguna reasigna dónde reside la capacidad representacional del modelo, porque la parametrización global del perceptrón multicapa no lo permite: añadir puntos de colocación en una región no dota a esa región de parámetros propios. La redistribución ocurre sobre el conjunto de evaluación, no sobre la representación.

Esta afirmación se refiere a las estrategias listadas en el cuadro \ref{tab:rejilla}, cuya parametrización es global. No debe extenderse sin matización a los trabajos de adaptación de grilla en arquitecturas de Kolmogorov-Arnold, que sí pueden reasignar elementos de representación. Estos trabajos se incorporan como una fila propia, mientras que el aporte que queda abierto para este trabajo es la asignación continua de centros y la asignación entre campos acoplados.

\section{Soporte local y asignación de capacidad}

Las redes de Kolmogorov-Arnold sustituyen las funciones de activación fijas sobre los nodos por funciones univariadas aprendibles sobre las aristas \cite{liu2024kan}. La construcción se apoya en el teorema de representación de Kolmogorov-Arnold, según el cual una función multivariada continua admite expresión como suma de composiciones de funciones univariadas.

\pendiente{Enunciar el teorema en forma cerrada. Esta es la única inyección ascendente del capítulo y se justifica porque la propiedad que se usa después depende de cómo se parametrizan las funciones univariadas.}

Cada función univariada se parametriza sobre un conjunto de elementos de base, y la propiedad relevante para este trabajo no es la familia concreta de base sino su \textbf{soporte}. Una base de soporte compacto satisface que modificar un coeficiente altera la salida únicamente en un entorno acotado del espacio de entrada. Bajo esa condición, la capacidad representacional del modelo queda asociada a regiones identificables del dominio, y no repartida de forma indistinguible sobre todos los parámetros.

Dos familias satisfacen la condición y se distinguen por cómo se ubican sus elementos:

\textbf{Bases de B-spline.} El soporte compacto es una consecuencia directa de la construcción por recursión sobre una partición del dominio. Los elementos están anclados a los nudos de esa partición, de modo que la ubicación de la capacidad es discreta: mover capacidad de una región a otra requiere reasignar nudos, insertarlos o refinar la grilla. Esta es la construcción del trabajo original sobre redes de Kolmogorov-Arnold \cite{liu2024kan} y la que emplean la mayoría de las variantes informadas por física.

\textbf{Bases radiales de soporte compacto.} Cada elemento es una función radial centrada en un punto del dominio y nula fuera de un radio finito. El soporte compacto se obtiene eligiendo un núcleo que se anula fuera de un radio —del tipo Wendland o \textit{bump}— y no una gaussiana, cuyo soporte es global y destruiría la propiedad de atribución que motiva toda esta sección. A diferencia del caso anterior, los centros no están anclados a una partición: son puntos del dominio que pueden tratarse como parámetros aprendibles [nuevo: Free-RBF-KAN, arXiv:2601.07760].

\pendiente{Incluir la definición formal del núcleo elegido y su radio, y una figura que contraste el soporte de una B-spline de orden k, una radial de soporte compacto y una gaussiana sobre el mismo intervalo.}

\textit{[Figura 6.x — Soporte de tres bases sobre el mismo intervalo: B-spline de orden $k$, radial de soporte compacto y gaussiana. La tercera se incluye para mostrar qué se pierde.]}

La distinción entre ambas familias no es de eficiencia sino de \textbf{granularidad de la asignación}. Con una base anclada a nudos, redistribuir capacidad es una operación discreta que ocurre entre etapas del entrenamiento. Con centros libres, la ubicación de la capacidad es una variable continua y diferenciable del mismo problema de optimización. Esa diferencia es la que hace posible que un mecanismo de focalización actúe sobre dónde reside la capacidad, y no solo sobre cómo se pondera el error donde la capacidad ya está.

La elección de una arquitectura informada por física basada en Kolmogorov-Arnold como sustrato de este trabajo se sigue de aquí, y no de una expectativa de superioridad frente a otras arquitecturas. Un mecanismo de focalización puede acoplarse a cualquier modelo informado por física que exponga una señal escalar sobre el dominio. En una arquitectura con parametrización global —donde cada peso influye sobre la salida en todo el dominio— el efecto de la focalización se confunde con la reorganización global de los parámetros y resulta difícil de atribuir. Con una base de soporte compacto, la atribución es directa: se puede observar si el mecanismo modificó la distribución espacial del aprendizaje sin que la observación quede mediada por ese efecto de confusión.

\pendiente{Decidir cuánta teoría de Kolmogorov-Arnold entra. Criterio de corte: entra lo necesario para entender por qué existe una señal interna y qué significa localmente. La discusión de interpretabilidad, poda simbólica y descubrimiento de fórmulas queda fuera.}

\section{Mecanismo de focalización adaptativa}

\subsection{Forma general — reformulada}

\textbf{Nota sobre el cambio.} La formulación anterior definía el mecanismo como una ponderación del residual, $\mathcal{L}_{\mathrm{foc}} = (1/N)\sum_i w(x_i;\theta)\lVert r(x_i;\theta)\rVert^2$ con $w=T[s(\cdot;\theta)]$. Esa forma pertenece a una familia cuyos miembros —esquemas de ponderación y de muestreo basados en el residual— han sido unificados formalmente como transformaciones convexas de un mismo funcional variacional [nuevo: Toscano et al., \textit{npj Artificial Intelligence}, 2026]. La reformulación que sigue traslada el mecanismo fuera de esa familia.

El mecanismo no modifica la contribución de cada punto del dominio a la función de pérdida. Modifica \textbf{dónde reside la capacidad representacional del modelo}, redistribuyendo los centros de la base durante el entrenamiento en función del estado interno del propio modelo.

Sea $\mu = \{\mu_1,\ldots,\mu_G\}$ el conjunto de centros de la base sobre una arista de la red. El mecanismo define un campo de densidad objetivo sobre el dominio a partir de una señal interna, y desplaza los centros hacia esa densidad.

\pendiente{Escribir la forma cerrada y elegir una variante antes del primer barrido: (a) deriva por gradiente, con centros entrenables y un término del objetivo que los atraiga hacia regiones de señal alta; o (b) reasignación por densidad, con reubicación periódica para igualar una densidad objetivo. La variante (a) sostiene el argumento de 6.5 y conserva la diferenciabilidad; la (b) es más fácil de comparar con la asignación de nudos, pero pierde la diferenciabilidad que motiva el cambio.}

La función de pérdida permanece en su forma estándar, ecuación \eqref{eq:perdida}. Esta es la diferencia estructural con las estrategias de la sección 6.4: aquellas modifican el funcional que se minimiza o la distribución sobre la que se lo estima; esta modifica el espacio de funciones sobre el que se busca la solución.

\subsection{Firma de tipos — revisada}

El mecanismo se define sobre tres objetos con dominio y codominio explícitos. La señal interna es un campo escalar sobre el dominio derivable del estado del modelo, $s : \Omega \to \mathbb{R}$. La transformación lleva ese campo a una densidad de asignación positiva y normalizada, $T : (\Omega \to \mathbb{R}) \to \mathcal{P}(\Omega)$. La densidad resultante gobierna la ubicación de los centros de la base.

Esta firma requiere dos condiciones del modelo sobre el que se instancia: que exponga una cantidad escalar resoluble sobre el dominio, y que su base tenga soporte compacto con elementos ubicables. La primera la satisface cualquier modelo informado por física; la segunda restringe la aplicación a arquitecturas del tipo descrito en la sección 6.5. El alcance de diseño y el alcance de validación se distinguen de forma explícita, y las conclusiones del trabajo se enuncian sobre el segundo.

\pendiente{Decidir si la firma admite además una componente sobre variables acopladas, T : (Omega x V -> R) -> P(Omega x V), donde V es el conjunto de campos físicos del sistema. Si se incluye, hay que introducirla desde aquí y no como extensión posterior, porque la sección de metodología la usa como segunda dimensión de la unidad de análisis.}

\subsection{Elección de la señal}

\pendiente{Desarrollar componente por componente, en orden de dependencia: extracción de la señal, normalización o transformación y punto de inyección. Cada uno debe incluir definición, justificación frente a la alternativa obvia y comportamiento en el caso degenerado. Restricción de diseño: la señal se define sobre el residual y la variación local de la solución, cantidades que cualquier modelo informado por física posee. Cualquier lectura desde los internos de la base spline es una instanciación opcional y se documenta como tal, no como parte de la definición.}

\subsection{Casos de reducción — revisados}

Si los centros permanecen fijos durante todo el entrenamiento, el modelo recupera la configuración base con grilla estática, que es la condición de control del diseño experimental.

Si la densidad objetivo es uniforme, el mecanismo reproduce el comportamiento de una grilla de resolución constante y la reasignación no transporta capacidad entre regiones.

Si la señal es el residual sin transformación y la reasignación se aplica de forma periódica en lugar de continua, el mecanismo se aproxima a los esquemas de asignación de nudos guiados por métricas de entrenamiento [nuevo: framework de funciones de densidad de importancia, arXiv:2601.18672]. Este límite es el más cercano de la literatura reciente y conviene enunciarlo explícitamente: la distancia del mecanismo propuesto respecto de él es la diferenciabilidad de la asignación y su extensión a la dimensión entre campos acoplados, no la idea de guiar la grilla por una señal.

Estos tres límites ubican la propuesta respecto de la literatura sin argumentación comparativa adicional.

\subsection{Costo — nota de alcance}

\pendiente{El análisis de complejidad cambia respecto de la formulación anterior. Ya no es el costo de calcular un campo de peso por época, sino el costo de evaluar la señal, construir la densidad y desplazar los centros, más el costo indirecto de que la base deje de estar sobre una partición fija. Incluir la complejidad adicional por época y la frecuencia de actualización.}

\section{Pruebas de Estres}
% =====================================================================

El mecanismo actúa sobre la distribución espacial del aprendizaje, de
modo que su evaluación exige problemas donde la dificultad se distribuya
de maneras distintas. Los problemas de este capítulo no se eligen por ser
los benchmarks habituales de la literatura, sino porque cada uno aísla un
tipo de no uniformidad,  ausente, estática en el espacio, móvil en el
tiempo, no localizada y repartida entre campos acoplados. El cuadro
\ref{tab:escalera} resume la correspondencia.

\begin{table}[htbp]
\centering
\caption{Problemas del banco y tipo de no uniformidad que aísla cada uno.}
\label{tab:escalera}
\begin{tabular}{lll}
\toprule
Problema & Tipo de no uniformidad & Función en el diseño \\
\midrule
Poisson oscilatorio & No localizada           & Verificación y control negativo \\
Burgers viscosa     & Estática en el espacio  & Localización espacial \\
Allen-Cahn          & Móvil en el tiempo      & Seguimiento de la región crítica \\
Navier--Stokes 2D   & Anclada a la geometría  & Localización con acoplamiento \\
Piezoeléctrico 2D   & Entre campos e interfaz & Acoplamiento físico \\
\bottomrule
\end{tabular}
\end{table}


\subsection{Problemas con referencia analítica}

La ecuación de Poisson establece la línea base y cumple además una
función de control. Con forzamiento suave, su solución no concentra
dificultad en ninguna región y la etapa se limita a verificar que la
implementación PIKAN reproduce soluciones conocidas dentro de la
tolerancia esperada. Con forzamiento oscilatorio de alta frecuencia, la
dificultad crece pero permanece repartida de forma homogénea sobre el
dominio, y el problema opera como control negativo: si el mecanismo
mejora el error también aquí, la mejora no proviene de la focalización
espacial y hay que buscar otra explicación, por ejemplo un efecto de
reponderación equivalente a un cambio de tasa de aprendizaje efectiva.
\pendiente{Fijar el dominio, la frecuencia del forzamiento y la tolerancia de aceptación; la tolerancia va en metodología.}

La ecuación de Burgers viscosa introduce el primer caso de dificultad
localizada. Para viscosidad pequeña, la solución desarrolla un gradiente
pronunciado en una región acotada del dominio mientras el resto permanece
suave \cite{raissi2019}. La localización es estática por lo que la región
crítica no se desplaza. Es el escenario más favorable al mecanismo y, por
esa misma razón, el que menos informa por sí solo.

\subsection{Localización móvil y localización estructural}

La ecuación de Allen-Cahn presenta una interfaz delgada que se desplaza
con el tiempo \cite{raissi2019}. La región de mayor dificultad existe
en cada instante pero cambia de posición, lo que somete al mecanismo a
una exigencia distinta de la anterior: no basta con concentrar capacidad
donde el residual fue alto, hay que reasignarla conforme la región se
mueve. Un mecanismo que concentre capacidad y no la libere degrada aquí
respecto de la configuración base, y ese comportamiento es observable.

Las ecuaciones de Navier--Stokes incorporan no linealidad en el término convectivo y acoplamiento entre el campo de velocidad y la presión. Para número de Reynolds moderado y régimen laminar, la dificultad se concentra enregiones determinadas por la geometría del dominio, como las capas límite
adyacentes a las paredes y la estela aguas abajo de un obstáculo
\cite{raissi2019,jin2021nsfnets}. La localización no se desplaza como en
Allen-Cahn ni es puntual como en Burgers, en donde  queda anclada a la frontera y
persiste durante todo el entrenamiento. Se evalúa si el mecanismo
distingue esas regiones sin que la señal quede dominada por el residual de
la ecuación de continuidad, cuyo comportamiento numérico difiere del de
la ecuación de momento.
\pendiente{Fijar la configuración y el número de Reynolds. El régimen turbulento y los flujos de alto Reynolds quedan fuera del alcance del trabajo y deben declararse como tales.}



\subsection{Sistema piezoeléctrico acoplado}

El efecto piezoeléctrico acopla el comportamiento mecánico y el eléctrico
de ciertos materiales. Las relaciones constitutivas lineales expresan el
tensor de esfuerzos y el desplazamiento dieléctrico en función de la
deformación y del campo eléctrico, y el sistema se cierra con las
ecuaciones de equilibrio mecánico y de Gauss
\cite{gonzalez2026pinnpiezo,lahmer2008}:
\begin{equation}
\nabla \cdot \sigma + \bm{f} = 0,
\qquad
\nabla \cdot \bm{D} = 0 .
\label{eq:piezo}
\end{equation}


\pendiente{escribir la forma matricial de las relaciones constitutivas y las relaciones cinemáticas, siguiendo la formulación de \cite{gonzalez2026pinnpiezo}. Enunciar, no derivar desde el tensor completo: la derivación no cambia ninguna interpretación de los resultados.}

La localización de la dificultad en
interfaces. En una viga bimorfa, las capas tienen polarización opuesta, y
la línea media donde se encuentran concentra la discrepancia entre la
solución aproximada y la de referencia \cite{gonzalez2026pinnpiezo}. La
segunda es la asimetría entre campos acoplados,  González et al. informan
errores relativos $L^2$ del orden de $10^{-1}$ para ambas componentes del
desplazamiento mecánico frente a $10^{-3}$ para el potencial eléctrico
\cite{gonzalez2026pinnpiezo}. La capacidad representacional se reparte de
forma desigual entre variables físicas acopladas, no solo entre regiones
del dominio.

\pendiente{decidir si el mecanismo se lee también sobre la dimensión entre variables o solo sobre la dimensión espacial. Si se incluye, la firma de tipos de la sección 6 debe admitirlo desde el principio.}

Conviene distinguir dos fuentes de error en estos mapas. La discrepancia
en la línea media corresponde a una interfaz física entre capas de
polarización opuesta. En cambio, la discrepancia en el extremo empotrado
aparece donde se impone una condición de Dirichlet de valor nulo y el
error relativo se amplifica cuando la solución de referencia se aproxima
a cero \cite{gonzalez2026pinnpiezo}. La primera constituye evidencia de
localización física; la segunda es, en parte, un artefacto de la métrica.
Esta distinción justifica reportar el residual físico junto con el error
relativo en el capítulo de resultados.

La referencia numérica para la validación proviene de una simulación por
elementos finitos con malla triangular bidimensional
\cite{gonzalez2026pinnpiezo,khan2018bimorph}, con los coeficientes de
material del PVDF diferenciados por capa según la orientación de
polarización.
\pendiente{trasladar el detalle de geometría, coeficientes y condiciones de frontera al capítulo de metodología; aquí solo la procedencia.}
